% project_report.tex — ERK-format report template for YOUR project (English).
% Fill every [PLACEHOLDER]. Same two-column style as the real ERK proceedings.
% Compile like any deck: pdflatex project_report.tex (erk.sty sits next to it).
\documentclass[a4paper]{article}

\usepackage[utf8]{inputenc}\usepackage[T1]{fontenc}
\usepackage{erk}
\usepackage{times}
\usepackage{graphicx}
\usepackage[english]{babel}
\usepackage[top=22.5mm, bottom=22.5mm, left=22.5mm, right=22.5mm]{geometry}
\usepackage{booktabs,hyperref}

% keep the ERK look: no footnotes on the cover page
\def\footnotemark{}

\begin{document}

\title{[Title of Your Project]}

\author{[Your Name]$^{1}$}

\affiliation{$^{1}$University of Ljubljana, Faculty of Computer and Information Science}

\email{Email: [your@student.uni-lj.si]}

\maketitle

\begin{abstract}{[English title of your paper]}
[English abstract, 200 words at most: what the software does, who it is for,
what the agents did, and the one result you can show. Write it last, readers
read it first.]
\end{abstract}

\section{Introduction}
What problem does your project solve, and who is the user? Why is it worth
building? Two short paragraphs maximum.

\section{Goals and objectives}
Numbered, checkable. For each goal someone can tell whether you met it.
\begin{itemize}
\item[G1.] [Goal one --- phrased so a reader can test it]
\item[G2.] [Goal two]
\end{itemize}

\section{Agentic stack used}
What you actually ran, pinned, with what each part did. Honest table, not
marketing. Model tags are part of the method --- pin them.

\begin{table}[h]
\centering\small
\begin{center}
\begin{tabular}{ | l | l | l | }
\hline 
  	extbf{Tool} & 	extbf{Model (pinned tag)} & 	extbf{What it did}\\
\hline
  pi 1.4.0 & dingo/qwen3.6:35b-a3b-mtp-q4\_K\_M & wrote, ran, fixed: what, when \\
  <tool 2> & <tag / version> & <its role> \\
\hline
\end{tabular}
\end{center}
\caption{The stack as used, not as advertised.}
\end{table}

\section{Method --- how you directed the agents}
How work was split, what files carried state (specs, progress files), how you
verified each agent output before committing. Your commit log is the
replication record: which prompt and which model produced which part.

\section{Results}
What works today; how you know (tests you ran, the demo, the counts). What
failed and stayed broken is a result too --- report it.

\section{Limitations and next steps}
What this version does not do; what you would do next; what you would change
about the process, not just the product.

\section{AI-use declaration}
What the agents produced versus what you produced; every model with its exact
pinned tag and date of use; what was auto-generated at release (model defaults).
This section is graded as part of the report; write it like a method section.

\small
\begin{thebibliography}{1}
\bibitem{repo} Repository: \url{[https://github.com/<username>-acs-project]} (public or instructor-accessible), branch/tag used for the demo.
\end{thebibliography}

\end{document}
